\FigureLayoutDeclare{img/2015/shanghai_liberal/q21_fig1.png}{4.8cm}{82bp 57bp 61bp 34bp}

\chapter{2015年上海卷（秋文）}

\section{填空题}

\begin{problem}
函数 \( f(x)=1-3\sin^{2}x \) 的最小正周期为\fillinblank{}.
\end{problem}

\begin{answer}
\(\pi\)
\end{answer}

\begin{problem}
设全集 \(U = \R\). 若集合 \( A = \{1, 2, 3, 4\} \)，\( B = \{x \mid 2 \leqslant x < 3\} \)，则 \( A \cap \complement_U B = \)\fillinblank{}.
\end{problem}

\begin{answer}
\(\{1,3,4\}\)
\end{answer}

\begin{problem}
若复数 \(z\) 满足 \( 3z + \overline{z} = 1 + \i \)，其中 \(\i\) 为虚数单位，则 \(z =\)\fillinblank{}.
\end{problem}

\begin{answer}
\(\frac{1}{4}+\frac{1}{2}\i\)
\end{answer}

\begin{problem}
设 \( f^{-1}(x) \) 为 \( f(x)=\frac{x}{2x+1} \) 的反函数，则 \( f^{-1}(2)= \)\fillinblank{}.
\end{problem}

\begin{answer}
\(-\frac{2}{3}\)
\end{answer}

\begin{problem}
若线性方程组的增广矩阵为 \( \begin{pmatrix} 2 & 3 & c_{1} \\ 0 & 1 & c_{2} \end{pmatrix} \)，解为 \( \begin{cases}
x = 3,\\
y = 5,
\end{cases} \) 则 \( c_{1}-c_{2}= \)\fillinblank{}.
\end{problem}

\begin{answer}
16
\end{answer}

\begin{problem}
若正三棱柱的所有棱长均为 \(a\)，且其体积为 \(16\sqrt{3}\)，则 \(a = \)\fillinblank{}.
\end{problem}

\begin{answer}
4
\end{answer}

\begin{problem}
抛物线 \( y^{2}=2px \) （\(p>0\)） 上的动点 \(Q\) 到焦点的距离的最小值为 1，则 \(p=\)\fillinblank{}.
\end{problem}

\begin{answer}
2
\end{answer}

\begin{problem}
方程 \( \log_{2}(9^{x-1}-5)=\log_{2}(3^{x-1}-2)+2 \) 的解为\fillinblank{}.
\end{problem}

\begin{answer}
2
\end{answer}

\begin{problem}
若 \(x\)，\(y\) 满足 \(\begin{cases}
x - y \geqslant 0,\\
x + y \leqslant 2,\\
y \geqslant 0,
\end{cases}\) 则目标函数 \(f = x + 2y\) 的最大值为\fillinblank{}.
\end{problem}

\begin{answer}
3
\end{answer}

\begin{problem}
在报名的 3 名男教师和 6 名女教师中，选取 5 人参加义务献血. 要求男、女教师都有，则不同的选取方式的种数为\fillinblank{}. （结果用数值表示）
\end{problem}

\begin{answer}
120
\end{answer}

\begin{problem}
在 \(\left(2x + \frac{1}{x^2}\right)^6\) 的二项展开式中，常数项等于\fillinblank{}. （结果用数值表示）
\end{problem}

\begin{answer}
240
\end{answer}

\begin{problem}
已知双曲线 \(C_1\)，\(C_2\) 的顶点重合，\(C_1\) 的方程为 \(\frac{x^2}{4} - y^2 = 1\)，若 \(C_2\) 的一条渐近线的斜率是 \(C_1\) 的一条渐近线的斜率的 2 倍，则 \(C_2\) 的方程为\fillinblank{}.
\end{problem}

\begin{answer}
\(\frac{x^2}{4}-\frac{y^2}{4}=1\)
\end{answer}

\begin{problem}
已知平面向量 \(\bs{a}\)，\(\bs{b}\)，\(\bs{c}\) 满足 \(\bs{a} \perp \bs{b}\)，且 \(\{|\bs{a}|, |\bs{b}|, |\bs{c}|\} = \{1, 2, 3\}\)，则 \(|\bs{a} + \bs{b} + \bs{c}|\) 的最大值是\fillinblank{}.
\end{problem}

\begin{answer}
\(3+\sqrt{5}\)
\end{answer}

\begin{problem}
已知函数 \( f(x) = \sin x \). 若存在 \(x_1\)，\(x_2\)，\(\dots\)，\(x_m\) 满足 \( 0 \leqslant x_1 < x_2 < \dots < x_m \leqslant 6\pi \)，且 \(|f(x_1) - f(x_2)| + |f(x_2) - f(x_3)| + \dots + |f(x_{m-1}) - f(x_m)| = 12\)（\(m \geqslant 2, m \in \N^*\)），则 \( m \) 的最小值为\fillinblank{}.
\end{problem}

\begin{answer}
8
\end{answer}

\section{单选题}

\begin{problem}
设 \(z_{1}\)，\(z_{2} \in \mathbf{C}\)，则“\(z_{1}\)，\(z_{2}\) 均为实数”是“\(z_{1} - z_{2}\) 是实数”的（\quad）

\choices
    {充分非必要条件}
    {必要非充分条件}
    {充要条件}
    {既非充分又非必要条件}
\end{problem}

\begin{answer}
A
\end{answer}

\begin{problem}
下列不等式中，与不等式 \(\frac{x + 8}{x^2 + 2x + 3} < 2\) 解集相同的是（\quad）

\choices
    {\((x + 8)(x^2 + 2x + 3) < 2\)}
    {\(x + 8 < 2(x^{2} + 2x + 3)\)}
    {\(\frac{1}{x^2 + 2x + 3} < \frac{2}{x + 8}\)}
    {\(\frac{x^2 + 2x + 3}{x + 8} > \frac{1}{2}\)}
\end{problem}

\begin{answer}
B
\end{answer}

\begin{problem}
已知点 \(A\) 的坐标为 \((4\sqrt{3}, 1)\)，将 \(OA\) 绕坐标原点 \(O\) 逆时针旋转 \(\frac{\pi}{3}\) 至 \(OB\)，则点 \(B\) 的纵坐标为（\quad）

\choices
    {\(\frac{3\sqrt{3}}{2}\)}
    {\(\frac{5\sqrt{3}}{2}\)}
    {\(\frac{11}{2}\)}
    {\(\frac{13}{2}\)}
\end{problem}

\begin{answer}
D
\end{answer}

\begin{problem}
设 \(P_{n}(x_{n}, y_{n})\) 是直线 \(2x - y = \frac{n}{n + 1}\) （\(n \in \N^*\)）与圆 \(x^{2} + y^{2} = 2\) 在第一象限的交点，则极限 \(\displaystyle\lim_{n \to \infty} \frac{y_{n} - 1}{x_{n} - 1} =\)（\quad）

\choices
    {\(-1\)}
    {\(-\frac{1}{2}\)}
    {\(1\)}
    {\(2\)}
\end{problem}

\begin{answer}
A
\end{answer}

\section{解答题}

\begin{problem}
如图，圆锥的顶点为 \(P\)，底面圆心为 \(O\)，底面的一条直径为 \(AB\)，\(C\) 为半圆弧 \(\widehat{AB}\) 的中点，\(E\) 为劣弧 \(\widehat{CB}\) 的中点. 已知 \(PO = 2\)，\(OA = 1\)，求三棱锥 \(P - AOC\) 的体积，并求异面直线 \(PA\) 与 \(OE\) 所成角的余弦值.

\bitmapfigure[float=inline,max width=0.42\linewidth]{img/2015/shanghai_liberal/q20_fig1.png}
\FloatBarrier
\end{problem}

\begin{solution}
在底面所在平面内建立直角坐标系，取 \(O\) 为原点，令 \(A=(-1,0,0)\)，\(B=(1,0,0)\)，\(C=(0,-1,0)\)，则 \(E=\left(\frac{\sqrt{2}}{2},-\frac{\sqrt{2}}{2},0\right)\). 由于 \(PO=2\)，可取 \(P=(0,0,2)\).

三角形 \(AOC\) 是直角三角形，且 \(OA=OC=1\)，所以其面积为 \(\frac{1}{2}\). 因而三棱锥 \(P-AOC\) 的体积为
\[
V=\frac{1}{3}\cdot\frac{1}{2}\cdot PO=\frac{1}{3}.
\]

直线 \(PA\) 的一个方向向量可取 \(\overrightarrow{AP}=(1,0,2)\)，直线 \(OE\) 的一个方向向量可取 \(\overrightarrow{OE}=\left(\frac{\sqrt{2}}{2},-\frac{\sqrt{2}}{2},0\right)\). 设两异面直线所成的角为 \(\theta\)，则
\[
\cos\theta=\frac{|\overrightarrow{AP}\cdot\overrightarrow{OE}|}{|\overrightarrow{AP}|\,|\overrightarrow{OE}|}=\frac{\frac{\sqrt{2}}{2}}{\sqrt{5}\cdot1}=\frac{\sqrt{10}}{10}.
\]
所以三棱锥 \(P-AOC\) 的体积为 \(\frac{1}{3}\)，异面直线 \(PA\) 与 \(OE\) 所成角的余弦值为 \(\frac{\sqrt{10}}{10}\).
\end{solution}

\begin{problem}
已知函数 \( f(x) = ax^2 + \frac{1}{x} \)，其中 \( a \) 为常数.
\begin{enumerate}
\item 根据 \(a\) 的不同取值，判断函数 \(f(x)\) 的奇偶性，并说明理由；
\item 若 \(a \in (1,3)\)，判断函数 \(f(x)\) 在 \([1,2]\) 上的单调性，并说明理由.
\end{enumerate}
\end{problem}

\begin{solution}
\begin{enumerate}
\item 函数的定义域为 \(\{x\in\R\mid x\neq0\}\)，关于原点对称. 对任意 \(x\neq0\)，有 \(f(-x)=ax^{2}-\frac{1}{x}\). 若 \(f\) 为偶函数，则应有 \(f(-x)=f(x)\)，从而 \(ax^{2}-\frac{1}{x}=ax^{2}+\frac{1}{x}\)，这对 \(x\neq0\) 不可能成立，所以 \(f\) 不是偶函数. 若 \(f\) 为奇函数，则应有 \(f(-x)=-f(x)\)，即 \(ax^{2}-\frac{1}{x}=-ax^{2}-\frac{1}{x}\)，于是 \(2ax^{2}=0\). 对任意 \(x\neq0\) 成立的充要条件是 \(a=0\). 因此，当 \(a=0\) 时，\(f(x)=\frac{1}{x}\) 为奇函数；当 \(a\neq0\) 时，\(f(x)\) 既不是奇函数，也不是偶函数.
\item 当 \(x\in[1,2]\) 且 \(a\in(1,3)\) 时，\(f\) 可导，且 \(f'(x)=2ax-\frac{1}{x^{2}}\). 因为 \(a>1\) 且 \(x\geqslant1\)，所以 \(f'(x)\geqslant2a-1>1>0\). 因此 \(f(x)\) 在 \([1,2]\) 上单调递增.
\end{enumerate}
\end{solution}

\begin{problem}
如图，\(O\)，\(P\)，\(Q\) 三地有直道相通，\(OP = 3\text{ 千米}\)，\(PQ = 4\text{ 千米}\)，\(OQ = 5\text{ 千米}\). 现甲，乙两警员同时从 \(O\) 地出发匀速前往 \(Q\) 地. 经过 \(t\) 小时，他们之间的距离为 \(f(t)\)（单位：千米）. 甲的路线是 \(OQ\)，速度为 \(5\text{ 千米/小时}\)，乙的路线是 \(OPQ\)，速度为 \(8\text{ 千米/小时}\). 乙到达 \(Q\) 地后在原地等待. 设 \(t = t_1\) 时，乙到达 \(P\) 地；\(t = t_2\) 时，乙到达 \(Q\) 地.
\begin{enumerate}
\item 求 \( t_{1} \) 与 \( f(t_{1}) \) 的值；
\item 已知警员的对讲机的有效通话距离是 3 千米. 当 \(t_1 \leqslant t \leqslant t_2\) 时，求 \(f(t)\) 的表达式，并判断 \(f(t)\) 在 \([t_1, t_2]\) 上的最大值是否超过 3？ 说明理由.
\end{enumerate}

\bitmapfigure[float=inline,max width=0.42\linewidth]{img/2015/shanghai_liberal/q21_fig1.png}
\FloatBarrier
\end{problem}

\begin{solution}
\begin{enumerate}
\item 乙沿 \(OP\) 以 \(8\text{ 千米/小时}\) 的速度行驶，所以 \(t_{1}=\frac{OP}{8}=\frac{3}{8}\). 取 \(O=(0,0)\)，\(P=(3,0)\)，\(Q=(3,4)\). 甲沿 \(OQ\) 行驶，经过 \(t_{1}\) 小时走过的路程为 \(5t_{1}=\frac{15}{8}\)，故此时甲的位置为 \(A=\left(\frac{9}{8},\frac{3}{2}\right)\). 因而
\[
f(t_{1})=AP=\sqrt{\left(3-\frac{9}{8}\right)^{2}+\left(0-\frac{3}{2}\right)^{2}}=\sqrt{\frac{225}{64}+\frac{144}{64}}=\frac{\sqrt{369}}{8}\text{ 千米}.
\]
\item 乙到达 \(Q\) 所需的时间为 \(t_{2}=\frac{OP+PQ}{8}=\frac{7}{8}\). 当 \(t_{1}\leqslant t\leqslant t_{2}\) 时，甲的位置为 \(M=(3t,4t)\)，乙已经经过 \(P\)，其位置为 \(N=(3,8t-3)\). 所以
\[
f(t)=MN=\sqrt{(3t-3)^{2}+(4t-8t+3)^{2}}=\sqrt{25t^{2}-42t+18},\qquad \frac{3}{8}\leqslant t\leqslant\frac{7}{8}.
\]
令 \(g(t)=f^{2}(t)=25t^{2}-42t+18\). 因为 \(g'(t)=50t-42\)，其顶点 \(t=\frac{21}{25}\) 位于区间 \(\left[\frac{3}{8},\frac{7}{8}\right]\) 内，且二次函数 \(g\) 开口向上，所以最大值在端点处取得. 计算得 \(g\left(\frac{3}{8}\right)=\frac{369}{64}\)，\(g\left(\frac{7}{8}\right)=\frac{25}{64}\). 因此
\[
\max_{t\in[t_{1},t_{2}]}f(t)=\frac{\sqrt{369}}{8}<3.
\]
所以 \(f(t)\) 在 \([t_{1},t_{2}]\) 上的最大值不超过 \(3\text{ 千米}\).
\end{enumerate}
\end{solution}

\begin{problem}
已知椭圆 \(x^{2} + 2y^{2} = 1\)，过原点的两条直线 \(l_{1}\) 和 \(l_{2}\) 分别与椭圆交于点 \(A\)，\(B\) 和 \(C\)，\(D\). 记 \(\triangle AOC\) 的面积为 \(S\).
\begin{enumerate}
\item 设 \(A(x_{1}, y_{1})\)，\(C(x_{2}, y_{2})\). 用 \(A\)，\(C\) 的坐标表示点 \(C\) 到直线 \(l_{1}\) 的距离，并证明 \(S = \frac{1}{2} |x_{1}y_{2} - x_{2}y_{1}|\)；
\item 设 \(l_{1}: y = kx\)，\(C\left(\frac{\sqrt{3}}{3}, \frac{\sqrt{3}}{3}\right)\)，\(S = \frac{1}{3}\). 求 \(k\) 的值.
\item 设 \(l_{1}\) 与 \(l_{2}\) 的斜率之积为 \(m\)，求 \(m\) 的值，使得无论 \(l_{1}\) 与 \(l_{2}\) 如何变动，面积 \(S\) 保持不变.
\end{enumerate}
\end{problem}

\begin{solution}
\begin{enumerate}
\item 直线 \(l_{1}\) 经过 \(O\) 和 \(A(x_{1},y_{1})\)，其方程可写为 \(y_{1}x-x_{1}y=0\). 因而点 \(C(x_{2},y_{2})\) 到直线 \(l_{1}\) 的距离为
\[
d=\frac{|y_{1}x_{2}-x_{1}y_{2}|}{\sqrt{x_{1}^{2}+y_{1}^{2}}}=\frac{|x_{1}y_{2}-x_{2}y_{1}|}{\sqrt{x_{1}^{2}+y_{1}^{2}}}.
\]
又 \(OA=\sqrt{x_{1}^{2}+y_{1}^{2}}\)，所以
\[
S=\frac{1}{2}OA\cdot d=\frac{1}{2}|x_{1}y_{2}-x_{2}y_{1}|.
\]
\item 因为 \(A\in l_{1}\)，且 \(l_{1}:y=kx\)，所以 \(y_{1}=kx_{1}\)，代入椭圆方程得 \(x_{1}^{2}(1+2k^{2})=1\)，从而 \(|x_{1}|=\frac{1}{\sqrt{1+2k^{2}}}\). 由 \(C\left(\frac{\sqrt{3}}{3},\frac{\sqrt{3}}{3}\right)\) 和第（1）问的面积公式，得
\[
\frac{2}{3}=2S=\frac{\sqrt{3}}{3}|x_{1}(1-k)|=\frac{\sqrt{3}}{3}\frac{|1-k|}{\sqrt{1+2k^{2}}}.
\]
两边平方并化简，得 \(3(1-k)^{2}=4(1+2k^{2})\)，即 \(5k^{2}+6k+1=0\). 因此 \(k=-1\) 或 \(k=-\frac{1}{5}\).
\item 设 \(l_{1}\) 的斜率为 \(k\)，则 \(l_{2}\) 的斜率为 \(\frac{m}{k}\). 由两条过原点的直线与椭圆的交点坐标，取绝对值后有
\[
S^{2}=\frac{\left(\frac{m}{k}-k\right)^{2}}{4(1+2k^{2})\left(1+2\left(\frac{m}{k}\right)^{2}\right)}=\frac{(m-k^{2})^{2}}{4(1+2k^{2})(k^{2}+2m^{2})}.
\]
令 \(u=k^{2}>0\). 若 \(m=0\)，则 \(S^{2}=\frac{u}{4(1+2u)}\)，随 \(u\) 变化，不能保持不变. 当 \(m\neq0\) 时，若 \(S\) 与 \(u\) 无关，则存在常数 \(c\)，使
\[
u^{2}-2mu+m^{2}=c\left(2u^{2}+(1+4m^{2})u+2m^{2}\right).
\]
比较二次项和常数项系数，得 \(c=\frac{1}{2}\). 再比较一次项系数，得 \(-2m=\frac{1+4m^{2}}{2}\)，即 \((2m+1)^{2}=0\)，所以 \(m=-\frac{1}{2}\). 反之，当 \(m=-\frac{1}{2}\) 时，
\[
S^{2}=\frac{\left(u+\frac{1}{2}\right)^{2}}{4(1+2u)\left(u+\frac{1}{2}\right)}=\frac{1}{8},
\]
与 \(u\) 无关. 因此所求 \(m=-\frac{1}{2}\).
\end{enumerate}
\end{solution}

\begin{problem}
已知数列 \(\{a_{n}\}\) 与 \(\{b_{n}\}\) 满足 \(a_{n + 1} - a_n = 2(b_{n + 1} - b_n)\)，\(n \in \N^*\).
\begin{enumerate}
\item 若 \( b_{n} = 3n + 5 \)，且 \( a_{1} = 1 \)，求 \( \{a_{n}\} \) 的通项公式；
\item 设 \(\{a_{n}\}\) 的第 \(n_0\) 项是最大项，即 \(a_{n_0} \geqslant a_n\)（\(n \in \N^*\)），求证： \(\{b_n\}\) 的第 \(n_0\) 项是最大项；
\item 设 \(a_{1} = 3\lambda < 0\)，\(b_{n} = \lambda^{n}\)（\(n \in \N^{*}\)），求 \(\lambda\) 的取值范围，使得对任意 \(m\)，\(n \in \N^{*}\)，\(a_{n} \neq 0\)，且 \(\frac{a_{m}}{a_{n}} \in \left(\frac{1}{6}, 6\right)\).
\end{enumerate}
\end{problem}

\begin{solution}
\begin{enumerate}
\item 因为 \(b_{n}=3n+5\)，所以 \(b_{n+1}-b_{n}=3\)，从而 \(a_{n+1}-a_{n}=6\). 数列 \(\{a_{n}\}\) 是首项为 \(a_{1}=1\)、公差为 \(6\) 的等差数列，因此
\[
a_{n}=1+6(n-1)=6n-5.
\]
\item 由已知关系得 \(a_{n+1}-2b_{n+1}=a_{n}-2b_{n}\)，所以 \(a_{n}-2b_{n}\) 是常数列，即 \(a_{n}=2b_{n}+a_{1}-2b_{1}\). 由 \(a_{n_{0}}\geqslant a_{n}\) 得
\[
2b_{n_{0}}+a_{1}-2b_{1}\geqslant2b_{n}+a_{1}-2b_{1},
\]
于是 \(b_{n_{0}}\geqslant b_{n}\). 因此 \(\{b_{n}\}\) 的第 \(n_{0}\) 项是最大项.
\item 由 \(a_{1}=3\lambda<0\) 得 \(\lambda<0\). 令 \(r=-\lambda>0\). 由恒等关系 \(a_{n}-2b_{n}=a_{1}-2b_{1}\)，得
\[
a_{n}=2\lambda^{n}+\lambda=\lambda+2\lambda^{n}.
\]
先求必要条件. 有 \(a_{1}=-3r\)，\(a_{2}=-r+2r^{2}=-r(1-2r)\). 因为任意两项之比都属于正数区间 \(\left(\frac{1}{6},6\right)\)，各项必须同号且不能为零，所以 \(0<r<\frac{1}{2}\). 又
\[
\frac{a_{1}}{a_{2}}=\frac{3}{1-2r}<6,
\]
解得 \(r<\frac{1}{4}\). 所以必要条件是 \(0<r<\frac{1}{4}\).

再证其充分性. 当 \(0<r<\frac{1}{4}\) 时，若 \(n\) 为奇数，则 \(a_{n}=-r-2r^{n}<0\)；若 \(n\) 为偶数，则 \(n\geqslant2\)，且 \(2r^{n}\leqslant2r^{2}<\frac{r}{2}\)，所以 \(a_{n}=-r+2r^{n}<0\). 因而所有 \(a_{n}\) 均非零且同为负数. 令 \(c_{n}=-a_{n}>0\). 当 \(n\) 为奇数时，\(c_{n}=r+2r^{n}\leqslant3r\)；当 \(n\) 为偶数时，\(c_{n}=r-2r^{n}\geqslant r-2r^{2}=r(1-2r)\). 因此对任意 \(m,n\in\N^*\)，有
\[
\frac{1-2r}{3}\leqslant\frac{a_{m}}{a_{n}}=\frac{c_{m}}{c_{n}}\leqslant\frac{3}{1-2r}.
\]
由于 \(0<r<\frac{1}{4}\)，有 \(\frac{1-2r}{3}>\frac{1}{6}\) 且 \(\frac{3}{1-2r}<6\). 所以题设条件成立. 综上，
\[
\lambda=-r\in\left(-\frac{1}{4},0\right).
\]
\end{enumerate}
\end{solution}
